%%%%%%%%%%%%%%%%%%%%%%%%%%%%%%%%%%%%%%%%%%%%%%%%%%%%%%%%%%%%%%%%%%%%%%%%

%%% LaTeX Template for AAMAS-2026 (based on sample-sigconf.tex)
%%% Prepared by the AAMAS-2026 Publication Chairs based on the version from AAMAS-2025. 

%%%%%%%%%%%%%%%%%%%%%%%%%%%%%%%%%%%%%%%%%%%%%%%%%%%%%%%%%%%%%%%%%%%%%%%%

\documentclass[sigconf,anonymous]{aamas} 

\usepackage{balance} 
\usepackage{algorithm}
\usepackage{algorithmic}
\usepackage{booktabs}
\usepackage{multirow}
\usepackage{colortbl}
\usepackage{xcolor}
\usepackage[section]{placeins}
\usepackage{amsmath}
\DeclareMathOperator*{\argmax}{arg\,max}
\DeclareMathOperator*{\argmin}{arg\,min}
\usepackage{caption}
\usepackage{subcaption}
\usepackage{todonotes}
\usepackage{orcidlink}
\usepackage{hyperref}

%%%%%%%%%%%%%%%%%%%%%%%%%%%%%%%%%%%%%%%%%%%%%%%%%%%%%%%%%%%%%%%%%%%%%%%%

\setcopyright{ifaamas}
\acmConference[AAMAS '26]{Proc.\@ of the 25th International Conference
on Autonomous Agents and Multiagent Systems (AAMAS 2026)}{May 25 -- 29, 2026}
{Paphos, Cyprus}{C.~Amato, L.~Dennis, V.~Mascardi, J.~Thangarajah (eds.)}
\copyrightyear{2026}
\acmYear{2026}
\acmDOI{}
\acmPrice{}
\acmISBN{}

%%%%%%%%%%%%%%%%%%%%%%%%%%%%%%%%%%%%%%%%%%%%%%%%%%%%%%%%%%%%%%%%%%%%%%%%

\acmSubmissionID{<<submission id>>}

\title[Counterfactual Gradient Alignment]{Counterfactual Gradient Alignment: Optimizing Directional Expert Supervision for Data-Efficient Learning}

\author{Arthur Pendragon}
\affiliation{
  \institution{Camelot Castle}
  \city{Camelot}
  \country{United Kingdom}}
\email{king.arthur@camelot.uk}

\author{Nimue}
\affiliation{
  \institution{The Lady's Lake}
  \city{Avalon}
  \country{United Kingdom}}
\email{lady.of.the.lake@avalon.uk}

\begin{abstract}
Typically, machine learning algorithms must find patterns in samples drawn from an unknown data distribution, with no prior contextual information about the input space or intuition about the problem they are solving. Humans, by contrast, can apply both intuition and prior knowledge to quickly solve new problems, requiring less training data. In this paper we investigate how to address this disparity by enriching traditional labels for supervised classification tasks; embedding additional information through gradient-based \emph{expert annotations}. We present a framework for \textbf{Counterfactual Gradient Alignment (CGA)} which (1) annotates existing observations with the direction in the feature space which indicates a change of class and (2) utilises a novel family of loss functions---ReLU, Softplus, and Sign-based variants---which reward negative model gradients along these directions. We demonstrate this approach across a diverse suite of benchmarks: synthetic 2D datasets (XOR, Circles, TwoMoons), image classification (MNIST), and multiple natural language sentiment classification tasks (IMDB, SST-2, SST-5). Our results reveal significant accuracy improvements in low-data regimes, such as a +7.17\% gain on IMDB and a +14.5\% improvement on complex 2D boundaries. Furthermore, we uncover a critical synergy between CGA and active learning, showing that directional supervision is uniquely effective when applied to informative boundary samples, whereas it can act as noise on random data subsets.
\end{abstract}

\keywords{Counterfactual Explanations, Gradient Supervision, Active Learning, Data Efficiency, Human-AI Interaction}

\newcommand{\BibTeX}{\rm B\kern-.05em{\sc i\kern-.025em b}\kern-.08em\TeX}

\begin{document}

\pagestyle{fancy}
\fancyhead{}

\maketitle 

\section{Introduction}
Highly tailored machine learning models typically require large amounts of data and may fail to generalise well to out-of-distribution examples, leading to unexpected or inaccurate predictions \cite{hand_classifier, Challen231, souza2020challenges}. Humans are, by contrast, relatively fast learners. We are able to utilise information outside of the usual data-label pair that is provided to machine learning systems, requiring fewer data points in order to draw conclusions and generalise to unseen data. It would be trivial, for example, for a human to adapt from film reviews to political tweets when performing sentiment classification, where a machine learning (ML) model may require refinement \cite{blitzer-etal-2007-biographies}. When operating in real environments the underlying distribution of observable data is often dynamic, subject to change and has limited availability \cite{Sahiner2023}. Particularly when errors carry a high cost, be it monetary or a risk to human safety, strong generalisation is a necessary component of any model which aims to assist or replace a highly adaptable human supervisor.

Whether we must maximise the utility of a few available samples or adapt to new observations quickly, the objective is often to learn the underlying distribution from a minimal number of examples while avoiding over-fitting. In online settings such as Active Learning, lack of data is addressed by the machine learning algorithm communicating with a (typically human) oracle who returns a new label for each query example, usually selected to reduce a form of uncertainty~\cite{prince2004does}. Formulations like Machine Teaching approach the problem by allowing humans to select minimal groups of samples that effectively describe concepts and patterns \cite{Zhu2015}. Both approaches can be said to improve data efficiency and the ability of models to generalise to new data through an informed selection of training data.

In this work we describe an alternative approach; rather than querying an oracle, we enrich the data in the form of more \emph{complex annotations}. We propose that embedding more information in the annotation and learning process can increase data efficiency, at the expense of annotation complexity. We test our hypothesis across multiple domains: (1) synthetic 2D datasets including XOR and Concentric Circles, (2) image classification on MNIST, and (3) sentiment classification on IMDB, SST-2, and SST-5. We present a method of enriching annotations with a direction vector which indicates the direction a sample must move in feature space in order to change the samples' classification; a \emph{counterfactual direction}. We also propose a novel family of loss functions which enforce negative model gradients in the direction of the counterfactual, with the intuition being that a change in class should happen along this direction. 

A central contribution of this work is the identification of a profound synergy between \textbf{Counterfactual Gradient Alignment (CGA)} and importance sampling. To ensure effective utilization of our limited annotation budget, we employ entropy-based importance sampling with a diversity parameter to strategically select the most informative samples for counterfactual annotation~\cite{he2014active}. This approach combines uncertainty-based selection with diversity constraints, maximizing the value of each human-provided counterfactual direction. Crucially, we investigate the synergy between these selection strategies and our alignment objectives, finding that the utility of directional supervision is profoundly dependent on sample quality.

\section{Related Work}
\label{related_work}

In this section we review state-of-the-art methods for including additional information in machine learning systems that go beyond data-label pairs, namely; human priors, counterfactual observations and gradient supervision.

\subsection{Embedding human priors in model training}
One significant motivation of this work is to enable human annotators to provide domain knowledge to models to improve data efficiency. \citeauthor{rieger2020interpretations} \cite{rieger2020interpretations} demonstrated the ability to embed human priors in their work developing contextual decomposition explanation penalisation (CDEP). CDEP enables embedding domain knowledge during model training to ignore spurious correlations, correct errors, and generalise to different types of dataset shifts. This method assumes that the explanation method is truthful and penalises model explanations where these do not align with explanations provided by an oracle. Domain knowledge can be encoded by a human or in an automated fashion. In a similar fashion, Ross et al. embed human priors as explanations and develop a loss function which balances cross entropy loss with an additional function to penalise misaligned model explanations \cite{ross2017right}. In their applications, they reinforce the magnitude of gradients to align to pre-supposed explanation; penalising the model where features are inappropriately neglected or considered to be too important.

\subsection{Training with counterfactuals}
Kaushik et al. \cite{Kaushik2020Learning} present a system for collecting counterfactually augmented data by tasking human annotators to minimally edit individual observations to produce a counterfactual on a subset of the IMDB dataset \cite{imdb}. They demonstrate that incorporating counterfactuals within the training data leads to improved generalisation and show that training purely with unaltered data or counterfactual pairs leads to over-fitting on spurious signals, while a combined dataset produces a model which is less sensitive to these signals.

In a follow up paper, Kaushik et al. investigate why counterfactually augmented data reduces spurious correlations \cite{kaushik2021explaining}. They propose that adding noise over causal features should worsen in-domain and out-of-domain performance, but that adding noise to non-causal features should improve relative performance on out-of-domain tasks. To test this hypothesis, they apply noise to the complement of the edited sections of text within their counterfactually augmented dataset. They repeat this process for sections of text identified by an attention-based classifier (BiLSTM with Self-Attention \cite{graves2005framewise}) and by gradient-based feature attribution \cite{li2016visualizing}, and find that human annotations served as the best augmentation for performance on out-of-distribution data. 

In our work we evaluate models trained on the counterfactually augmented data produced both by humans and algorithmically. However, we take a different approach to learning from these counterfactual augmentations through gradient supervision.

\subsection{Learning from counterfactuals through gradient supervision}
\label{subsec:gradient_supervision}
Teney et al. \cite{teney2020learning} introduce a novel training objective called ``gradient supervision", which involves penalising the angle between the gradients of a neural network classifier and the vector difference between counterfactual examples in the input space. Using a counterfactual, Teney et al. define a loss function which aims to minimise the angle between the model gradient and vector pointing to the counterfactual as 
\begin{equation}
        \mathcal{L}_{GS} \left(\mathbf{g}_i,\hat{\mathbf{g}}_i \right)  = 
        1 - \langle\mathbf{g}_i,\hat{\mathbf{g}}_i \rangle /
        \left( \lVert \mathbf{g}_i \rVert \lVert \hat{\mathbf{g}}_i \rVert \right)
        \label{eq:gradient_supervision}
\end{equation}
where $\langle \cdot, \cdot \rangle$ is the inner product between two vectors, $\mathbf{g}_i$ denotes the gradient of the model with respect to an observation $\mathbf{x}_i$ and $\hat{\mathbf{g}}_i$ is the vector between counterfactual examples.

Eq.\ref{eq:gradient_supervision} is used as a regulariser alongside the conventional cross-entropy loss in order to align the model's decision surface between pairs of counterfactual examples. The authors demonstrate the effectiveness of their approach on tasks known for their susceptibility to poor generalisation due to biases in the training data; visual question answering, multi-label image classification, sentiment analysis, and natural language inference. Where Teney et al. assume alignment between the feature space and the output space and attempt to minimise the angle between the gradient of the model and the counterfactual direction, our proposed loss functions focus on the \emph{directional derivative}, exploring smoother landscapes (Softplus) and bounded signals (Sign) to achieve more robust decision boundaries.

\subsection{Importance Sampling}
Importance sampling techniques have been widely adopted in active learning to improve annotation efficiency by prioritizing samples that contribute most to model learning. Early work by Settles and Craven~\cite{settles2008active} introduced cost-sensitive active learning, demonstrating that accounting for variable annotation costs can significantly improve learning curves compared to random sampling. He et al.~\cite{he2014active} extended this idea by combining uncertainty sampling with representativeness through a hybrid strategy leveraging kernel k-means clustering. Their approach selects samples based on a weighted combination of prediction entropy and similarity to the unlabeled pool, ensuring both informativeness and diversity in each batch.

Doucet et al.~\cite{doucet2024bridging} bridged uncertainty and diversity by introducing a diversity parameter into entropy-based sampling, striking a balance between selecting highly uncertain samples and maintaining broad coverage of the feature space. Our work demonstrates that the benefits of gradient supervision are uniquely coupled with these sampling techniques, flourishing on boundary samples while being potentially harmful on random data subsets.

\section{Gradient-Based Learning}
In traditional supervised machine learning for classification, we assume a dataset of $N$ random independent identically distributed (i.i.d.) observations {$\{\mathbf{x_i}, y_i\}$} $\sim P(\mathbf{x},y)$, where $\mathbf{x_i}\in\mathcal{R}^d$ and $y_i$ is the target class label~\cite{hastie2009elements}. Here, let us consider instead that we have access to triplets of the form of {$\{\mathbf{x_i}, y_i, \mathbf{d_i}\}$}, where $\mathbf{d_i}\in\mathcal{R}^d$ defines a unit vector from a point $\mathbf{x_i}$ which points towards a proximal example of the opposing class; a \emph{counterfactual}.

Formally, a counterfactual $\hat{\mathbf{x}}_i$ is defined as the closest point to $\mathbf{x_i}$ for which the class is not the same as $y_i$ according to model $g(\mathbf{x})$ and a distance function $D$;
\begin{equation}\label{eq:counterfactual}
\begin{aligned}
    \hat{\mathbf{x}}_i =& \argmin_{\mathbf{x}} D(\mathbf{x_i}-\mathbf{x}) \ \ \ \forall \ \mathbf{x} \\ &	\text{such that} \ y_i \neq g(\mathbf{x})
\end{aligned}
\end{equation}
where $g(\mathbf{x})$ can define a model to be queried, or in our case an expert human annotator. We propose that the direction from $\mathbf{x_i}$ towards a local counterfactual should indicate a decrease in prediction probabilities generated from our classifier $f(\mathbf{x})$ in that direction. We can define the counterfactual direction as a unit vector given by; 
\begin{equation}
    \mathbf{d_i}=\frac{\mathbf{\hat{x}}_i - \mathbf{x_i}}{||\mathbf{\hat{x}}_i - \mathbf{x_i}||}. \label{eq:dir_unit}
\end{equation}
The \textbf{Directional Derivative} $d$ is computed as the projection of the gradient of the model's score for the true class onto this vector:
\begin{equation}
    d = \nabla_x f_{y}(x) \cdot \mathbf{d}_{i}
\end{equation}
A negative $d$ indicates desired behavior (score drops towards the counterfactual). To align the gradients of our model with the direction $\mathbf{d_i}$, we propose and evaluate three loss variants:

\begin{enumerate}
    \item \textbf{ReLU Loss}: $\mathcal{L}_{\text{ReLU}} = \max(0, d)$. This applies a linear penalty only if the derivative is positive (confidence is increasing towards the counterfactual).
    \item \textbf{Softplus Loss}: $\mathcal{L}_{\text{Softplus}} = \log(1 + e^{\beta d})$. A smooth version of ReLU that encourages a steep confidence drop-off even when the model is already \emph{correct}.
    
    \item \textbf{Sign Loss}: $\mathcal{L}_{\text{Sign}} = \tanh(\beta d) + 1$. This focuses purely on the sign of the derivative, providing a bounded signal that is robust to outliers.
\end{enumerate}

Unlike standard cross entropy loss $\mathcal{L}_{\textrm{CE}}$, which is calculated using the ground truth class labels, our novel direction loss signal $\mathcal{L}_{d}$ requires that at least some of the observations in our dataset contain counterfactual directions. We therefore incorporate both loss terms into our loss function, introducing a control parameter $\alpha \in [0,1]$, which modifies the ratio of standard cross entropy loss to our direction loss which is used in training. Equation~\ref{eq:lambda} defines our total loss, $\mathcal{L}$.

\begin{align}\label{eq:lambda}
    \mathcal{L} = (1-\alpha)\mathcal{L}_{\textrm{CE}} + \alpha \mathcal{L}_{d}
\end{align}

\section{Experimental Setup}
Evaluation of our method is applicable to datasets where counterfactuals are available within the training data. We utilise several open-source datasets where counterfactuals are provided by humans \cite{Kaushik2020Learning} or generated automatically. The objective of our experiments is to measure how incorporating counterfactuals using Eq.\ref{eq:lambda} improves data efficiency. To this end we utilise small subsets of the available training data to simulate data scarcity.

\subsection{2D Synthetic Tasks}
We evaluate CGA on three synthetic benchmarks: XOR, Concentric Circles, and Two Moons. XOR features four clusters where adjacent clusters have different labels. Circles features two concentric distributions. Two Moons features two interleaved semi-circles. These tasks allow for clear visualization of decision boundary shifts. We test dataset sizes of 10, 30, and 100 samples. The model is a simple 2-layer MLP.

\subsection{Image Classification (MNIST)}
We use the MNIST dataset with three types of counterfactual paths:

\begin{itemize}
    \item \textbf{FACE}: Plausible paths generated using the Feasible Path algorithm \cite{poyiadzi2020face}.
    \item \textbf{PCA}: Shortest paths in PCA-transformed feature space.
    \item \textbf{Raw}: Euclidean paths in raw pixel space.
\end{itemize}

The model is an optimized CNN (OptimizedMNISTConv) with two convolutional layers and a 512-dimensional dense head. We tested training set sizes of 10, 50, and 200 samples.

\subsection{Text Classification}
We utilize three sentiment classification benchmarks: IMDB (binary), SST-2 (binary), and SST-5 (5-class). Reviews are preprocessed (HTML removed, tokenized, stemmed) and represented as average-pooled 50-dimensional embeddings using a Bag-of-Words neural network. We simulate scarcity with subsets ranging from 10 to 200 samples.

\subsection{Active Learning Integration}
To test the dependency on sample quality, we compared two data acquisition strategies: (1) \textbf{AL ON}: Entropy-based importance sampling with a diversity parameter to select boundary samples; (2) \textbf{AL OFF}: Purely random selection. Training was performed for 5 epochs using the AdamW optimizer.

\section{Results and Analysis}


\subsection{Synthetic 2D Benchmarks}
We evaluate CGA's ability to regularize complex topologies was most evident in the Concentric Circles task, where directional information guided the model to the correct boundary where standard CE failed to generalize from sparse points. We also evaluated performance on XOR and Two Moons benchmarks. While these tasks are more easily separable by standard objectives, CGA maintained competitive performance, providing marginal gains in the ultra-low data XOR regime (Size 10). The results across these synthetic topologies are summarized in Table \ref{tab:2d_results}.

\begin{table}[h]
\centering
\caption{Validation Accuracy on Synthetic 2D Topologies (Size 100 for Circles, Size 10 for XOR/Moons).}
\label{tab:2d_results}
\begin{tabular}{@{}lcccc@{}}
\toprule
\textbf{Dataset} & \textbf{Baseline (CE)} & \textbf{Best CGA} & \textbf{Variant} & \textbf{Gain} \\ \midrule
Circles & 57.50\% & \textbf{72.00\%} & Sign & \textbf{+14.50\%} \\
XOR     & 74.50\% & \textbf{75.50\%} & Sign & +1.00\% \\
TwoMoons & \textbf{81.50\%} & 81.00\% & Softplus & -0.50\% \\ \bottomrule
\end{tabular}
\end{table}


\subsection{NLP Benchmarks and AL Synergy}
As shown in Table \ref{tab:nlp_results}, CGA achieves substantial gains on IMDB and SST-5 sentiment tasks when paired with Active Learning (AL). The Softplus variant was the most consistent performer, providing a +7.17\% gain on IMDB.

\begin{table}[h]
\centering
\caption{Validation Accuracy on IMDB and SST-5 benchmarks (AL ON, 5 Epochs).}
\label{tab:nlp_results}
\begin{tabular}{@{}lcccc@{}}
\toprule
\textbf{Dataset} & \textbf{Size} & \textbf{Baseline (CE)} & \textbf{Best CGA} & \textbf{Gain} \\ \midrule
IMDB & 200 & 55.94\% & \textbf{63.11\% (SP)} & \textbf{+7.17\%}
 \\IMDB & 100 & 51.43\% & \textbf{56.56\% (Sign)} & \textbf{+5.13\%}
 \\SST-5 & 100 & 53.06\% & \textbf{56.46\% (ReLU)} & \textbf{+3.40\%}
 \\ \bottomrule
\end{tabular}
\end{table}

Crucially, CGA's performance is profoundly sensitive to sample informativeness. Without AL (Table \ref{tab:al_synergy}), the directional signal degrades performance, acting as noise on random training subsets.

\begin{table}[h]
\centering
\caption{Effect of Active Learning (AL) on CGA Utility (Gain over CE Baseline at Size 100).}
\label{tab:al_synergy}
\begin{tabular}{lcc}
\toprule
\textbf{Dataset} & \textbf{AL Enabled (ON)} & \textbf{AL Disabled (OFF)} \\ \midrule
Gain (IMDB) & \textbf{+5.13\%} & -2.26\%
 \\Gain (SST-2) & \textbf{+1.36\%} & -2.72\%
 \\Gain (SST-5) & \textbf{+3.40\%} & -4.77\%
 \\ \bottomrule
\end{tabular}
\end{table}

\subsection{Image Classification and Path Feasibility}
We evaluated the impact of directional quality on MNIST. FACE paths, which prioritize plausible data transitions, provided superior results at larger sizes compared to raw geometric paths (Table \ref{tab:mnist_paths}).

\begin{table}[h]
\centering
\caption{MNIST Accuracy Comparison across Path Types (AL ON, Size 200).}
\label{tab:mnist_paths}
\begin{tabular}{lcccc}
\toprule
\textbf{Path Type} & \textbf{Baseline} & \textbf{Best CGA} & \textbf{Variant} & \textbf{Gain} \\ \midrule
Raw & 83.50\% & 84.52\% & Sign & +1.02\%
 \\PCA & 83.50\% & 84.56\% & Softplus & +1.06\%
 \\FACE & 83.50\% & \textbf{84.78\%} & Softplus & \textbf{+1.28\%}
 \\ \bottomrule
\end{tabular}
\end{table}

\section{Discussion}
The consistent superiority of the \textbf{Softplus} variant suggests that smooth, persistent gradient alignment is more effective than hard hinge penalties. Softplus ensures the model maintains a high-margin confidence drop even when \"correct.\" The failure of CGA on random data subsets reveals that directional supervision is a high-precision signal. For \"easy\" samples, the counterfactual direction may align with irrelevant features or conflict with the classification signal. However, for boundary samples, the directional signal provides critical structural information. 

Finally, the MNIST path analysis indicates that \textbf{feasibility matters}. Directions that follow the data manifold (FACE) provide better regularizing signals than raw Euclidean vectors, especially as the data size increases and the manifold becomes more complex.

\section{Conclusions}
We present Counterfactual Gradient Alignment (CGA), a method for teaching models the local geometry of decision boundaries using expert-provided directions. We evaluated three novel loss variants, demonstrating significant improvements in data-scarce vision and NLP tasks. Crucially, we identified a critical synergy between CGA and active learning, proving that directional supervision thrived on informative boundary samples. Future work will investigate human-in-the-loop interfaces for real-time acquisition of these machine-interpretable signals.

\balance

\bibliographystyle{ACM-Reference-Format}
\bibliography{sample}

\appendix

\section{Entropy-based importance sampling with diversity}
Where gradients have been used previously as a natural method for providing human-interpretable explanations~\cite{simonyan2013deep}, this paper aims to reverse this pipeline and provide machine-interpretable supervision.

To ensure that our framework selects the most informative samples for counterfactual annotation while maintaining diversity across the feature space, we employ an entropy-based importance sampling strategy with an integrated diversity parameter. This approach addresses the common limitations of purely uncertainty-based or diversity-based sampling methods by combining their respective strengths~\cite{he2014active,doucet2024bridging}.

\textbf{Implementation Details.} Our sampling procedure operates iteratively during training. At each selection round, we first compute the softmax probabilities from the current model's logits: $\mathbf{p}_i = \text{softmax}(\mathbf{f}_\theta(\mathbf{x}_i))$. The uncertainty of each sample is then quantified using the entropy of these probability distributions: $U(\mathbf{x}_i) = -\sum_{y \in \mathcal{Y}} p_{i,y} \log p_{i,y}$, where $p_{i,y}$ denotes the predicted probability for class $y$. Samples with high entropy indicate regions of significant uncertainty.

However, relying solely on uncertainty sampling can lead to redundant selection near decision boundaries. To address this, we incorporate a diversity parameter $\beta$ that encourages the selection of samples spanning different regions. We use the learned embeddings $\mathbf{e}_i$ from our model's intermediate representations to compute representativeness $Rep(\mathbf{x}_i)$ based on similarity to other unlabeled instances: $Rep(\mathbf{x}_i) = \frac{1}{|U|} \sum_{\mathbf{x}_j \in U} \exp\left(-\frac{||\mathbf{e}_i - \mathbf{e}_j||^2}{2\sigma^2}\right)$. The combined information content is then: $Info(\mathbf{x}_i) = U(\mathbf{x}_i)^\alpha \cdot Rep(\mathbf{x}_i)^\beta$.

This iterative process ensures that our model continuously improves its understanding of challenging regions while maintaining global structural coverage. For our gradient supervision framework, this is particularly important, as diverse counterfactual directions provide richer supervision signals that improve model generalization across the entire feature space.

\section{Shortest path definitions}
\paragraph{Raw Distance Counterfactuals}
Generates counterfactuals by finding the nearest instance in the target class using standard Euclidean distance: $\mathbf{x}_{cf} = \arg\min_{\mathbf{x} \in \mathcal{X}_{target}} ||\mathbf{x} - \mathbf{x}_0||_2$. This approach provides the most direct path but may traverse low-density regions.

\paragraph{PCA Distance Counterfactuals}
Addresses feature correlation by transforming data into principal component space before finding the nearest counterfactual: $\mathbf{z}_{cf} = \arg\min_{\mathbf{z} \in \mathcal{Z}_{target}} ||\mathbf{z} - \mathbf{z}_0||_2$, where $\mathbf{z} = \mathbf{P}^T\mathbf{x}$. This can reveal semantically meaningful proximity relationships obscured by raw feature scales.

\end{document}
%%%%%%%%%%%%%%%%%%%%%%%%%%%%%%%%%%%%%%%%%%%%%%%%%%%%%%%%%%%%%%%%%%%%%%%%
